\documentclass[12pt]{article}
\usepackage[utf8]{inputenc}
\usepackage[T1]{fontenc}
\usepackage{hyperref}
\title{Sochiapam Chinantec Language Resources for K-12 Curriculum Development}
\author{Compiled by Claude Sonnet 5.5 (claude-sonnet-5-5) for Andr\'as Kornai}
\date{}
\begin{document}
\maketitle

\section{Introduction}
Sochiapam Chinantec (ISO 639-3 \texttt{cso}) is an Otomanguean language of San Pedro Sochiapam, Oaxaca, Mexico, with about 3,600 speakers in older figures. It has seven tones.

\section{Primary Educational Resources}
No primers, textbooks or school programmes were found in the search results.

\section{Linguistic and Reference Materials}
\begin{itemize}
\item Sochiapam Chinantec Living Dictionary, with audio, produced with the Living Tongues Institute, the Municipalidad de San Pedro Sochiapam and the Observatorio de Lenguas Indígenas: \url{https://livingdictionaries.app/sochiapam-chinantec}
\item David Foris, \emph{A grammar of Sochiapan Chinantec} (SIL, Dallas, 2000). It is cited in the ASJP entry; no direct link to the grammar itself was found. ASJP language page: \url{https://asjp.clld.org/languages/CHINANTEC_SOCHIAPAN}
\item Numeral data page, Max Planck Institute for Evolutionary Anthropology: \url{https://lingweb.eva.mpg.de/channumerals/Chinantec-Sochiapam.htm}
\item Wikipedia ISO 639 entry: \url{https://wikipedia.com/wiki/ISO_639:cso}
\end{itemize}

\section{Religious and Community Resources}
\begin{itemize}
\item Global Recordings Network, gospel audio: \url{https://www.globalrecordings.net/en/language/cso}
\end{itemize}

\section{Digital and Audio Resources}
\begin{itemize}
\item Living Dictionary (above), with recorded audio.
\item Wikitongues language page: \url{https://wikitongues.org/languages/cso/}
\end{itemize}

\section{Curriculum Development Context}
The Living Dictionary is the main reusable resource and could seed vocabulary material. The search results did not show any school-oriented material or a downloadable wordlist.

\end{document}